\subsection{有限维单 g-模的权}

命题 1. 设 V 为有限维 g-模。

\begin{enumerate}
\def\labelenumi{(\roman{enumi})}
\item
  V 的所有权都属于 P。
\item
  \(V = \bigoplus_{\mu \in P} V^{\mu}.\)
\item
对于所有 \(\mu \in \mathfrak{h}^*\)，\(\mathrm{V}^{\mu}\) 是这样的 \(x \in \mathrm{V}\) 的集合：满足 \(h.x = \mu(h)x\)（对所有 \(h \in \mathfrak{h}\)）。
\end{enumerate}

对于所有 \(\alpha \in \mathrm{B}\)，存在一个从 \(\mathfrak{sl}(2,k)\) 到 \(\mathfrak{g}\) 的同态，将 \(H\) 映到 \(H_{\alpha}\)。因此，根据 §1，第 2 节，命题 2 的推论，\((H_{\alpha})_{\mathrm{V}}\) 可对角化，且其特征值为整数。于是，\((H_{\alpha})_{\mathrm{V}}\) 所成的集合（其中 \(\alpha \in \mathrm{B}\)）是可对角化的（《代数》，第 VII 章，§5，第 6 节，命题 13）。因此，对于所有 \(h \in \mathfrak{h}\)，\(h_{\mathrm{V}}\) 都可对角化。根据第 VII 章，§1，第 3 节，命题 9，\(\mathrm{V} = \bigoplus_{\mu \in \mathfrak{h}^*} \mathrm{V}^{\mu}\)。另一方面，如果 \(\mathrm{V}^{\mu} \neq 0\)，则前述结果表明 \(\mu(H_{\alpha}) \in \mathbf{Z}\)，且这对所有 \(\alpha \in \mathrm{B}\) 都成立，所以 \(\mu \in \mathrm{P}\)。这就证明了 (i) 和 (ii)。同理可见 \(\mathfrak{h}_{\mathrm{V}}\) 可对角化，因而得到 (iii)。

推论。设 \(\rho\) 是 \(\mathfrak{g}\) 的有限维表示，\(\Phi\) 是与 \(\rho\) 相关的双线性型。

\begin{enumerate}
\def\labelenumi{(\roman{enumi})}
\item
  若 \(x, y \in \mathfrak{h}_{\mathbf{Q}}\)，则 \(\Phi(x, y) \in \mathbf{Q}\)，且 \(\Phi(x, x) \in \mathbf{Q}_{+}\)。
\item
  若 \(\rho\) 是单射，则 \(\Phi\) 在 \(\mathfrak{h}\) 上的限制是非退化的。
\end{enumerate}

断言 (i) 由命题 1 得出，因为 P 的元素在 \(\mathfrak{h}_{\mathbf{Q}}\) 上取有理值。如果 \(\rho\) 是单射，则 \(\Phi\) 是非退化的（第 I 章，§6，第 1 节，命题 1），所以 \(\Phi\) 在 \(\mathfrak{h}\) 上的限制是非退化的（第 VII 章，§1，第 3 节，命题 10 (iii)）。

引理 1. 设 \(\mathrm{V}\) 是一个 \(\mathfrak{g}\)-模，\(\rho\) 是相应的 \(\mathfrak{g}\) 表示。

\begin{enumerate}
\def\labelenumi{(\roman{enumi})}
\tightlist
\item
  若 \(a\) 是 \(\mathfrak{g}\) 的幂零元，且 \(\rho(a)\) 是局部幂零的，
\end{enumerate}

\[
\rho (e ^ {\mathrm{ad} a} b) = e ^ {\rho (a)} \rho (b) e ^ {- \rho (a)}
\]

对所有 \(b \in \mathfrak{g}\) 都成立。

\begin{enumerate}
\def\labelenumi{(\roman{enumi})}
\setcounter{enumi}{1}
\tightlist
\item
  若 \(\alpha \in \mathrm{R}\)，且经 \(\rho\) 作用，\(\mathfrak{g}^{\alpha}\) 和 \(\mathfrak{g}^{-\alpha}\) 中的元素的像都是局部幂零的，则 \(V\) 的权集在反射 \(s_{\alpha}\) 下稳定。
\end{enumerate}

在 (i) 的假设下，有 \(\rho((\mathrm{ad}a)^n b) = (\mathrm{ad}\rho(a))^n\rho(b)\) 对所有 \(n\geq 0\) 成立，所以 \(\rho(e^{\mathrm{ad}a}b) = e^{\mathrm{ad}\rho(a)}\rho(b)\)。另一方面，

\[
e ^ {\mathrm{ad} \rho (a)} \rho (b) = e ^ {\rho (a)} \rho (b) e ^ {- \rho (a)}
\]

这就是第 VII 章，§3，第 1 节，引理 1 的断言 (ii)。

现在采用 (ii) 的假设。令 \(\theta_{\alpha} = e^{\mathrm{ad}X_{\alpha}}e^{\mathrm{ad}X_{-\alpha}}e^{\mathrm{ad}X_{\alpha}}\)。由 (i)，存在 \(S\in \mathbf{GL}(V)\)，使得 \(\rho (\theta_{\alpha}b) = S\rho (b)S^{-1}\) 对所有 \(b\in \mathfrak{g}\) 成立。现在 \(\theta_{\alpha}|\mathfrak{h}\) 是 \(s_{\alpha}\) 的转置（§2，第 2 节，引理 1）。设 \(\lambda\) 是 V 的一个权。存在 V 的非零元素 \(x\)，使得 \(\rho (h)x = \lambda (h)x\) 对所有 \(h\in \mathfrak{h}\) 成立。于是

\[
\rho (h) \mathrm{S} ^ {- 1} x = \mathrm{S} ^ {- 1} \rho (^ {t} s _ {\alpha} h) x = \mathrm{S} ^ {- 1} \lambda (^ {t} s _ {\alpha} h) x = (s _ {\alpha} \lambda) (h) \mathrm{S} ^ {- 1} x
\]

对所有 \(h \in \mathfrak{h}\) 都成立。因此，\(s_{\alpha}\lambda\) 是 V 的一个权。

命题 2. 设 \(\mathrm{V}\) 是有限维 \(\mathfrak{g}\)-模，且 \(s \in \operatorname{Aut}_0(\mathfrak{g})\)。

\begin{enumerate}
\def\labelenumi{(\roman{enumi})}
\item
  存在 \(\mathrm{S} \in \mathbf{GL}(\mathrm{V})\)，使得 \((s(x))_{\mathrm{V}} = \mathrm{S}x_{\mathrm{V}}\mathrm{S}^{-1}\) 对所有 \(x \in \mathfrak{g}\) 成立。
\item
  若 \(s \in \mathrm{Aut}_e(\mathfrak{g})\)，则可将 \(S\) 取为 \(\mathbf{SL}(V)\) 的元素，使所有 \(\mathfrak{g}\)-子模在 \(V\) 中保持不变。
\end{enumerate}

断言 (ii) 由引理 1 (i) 得出。现在设 \(s \in \mathrm{Aut}_0(\mathfrak{g})\)，并以 \(\rho\) 表示由 V 定义的 \(\mathfrak{g}\) 表示。由 (ii)，表示 \(\rho\) 与 \(\rho \circ s\) 在扩张标量后变得等价。因此它们本来就等价（第 I 章，§3，第 8 节，命题 13），从而 S 存在。

备注 1. 设 S 满足命题 2 (i) 中的条件，并令 \(\mathfrak{h}' = s(\mathfrak{h})\)；以 \(s^*\) 表示同构 \(\lambda \mapsto \lambda \circ s^{-1}\)，它从 \(\mathfrak{h}^*\) 到 \(\mathfrak{h}'^*\)。显然有

\[
\mathrm{S} (\mathrm{V} ^ {\lambda}) = \mathrm{V} ^ {s ^ {*} \lambda}.
\]

特别地：

推论 1. 同构 \(s^*\) 将 V 关于 \(\mathfrak{h}\) 的权映到 V 关于 \(\mathfrak{h}'\) 的权；对应的权具有相同的重数。

推论 2. 设 \(w \in \mathrm{W}\)。对于所有 \(\lambda \in \mathfrak{h}^*\)，向量子空间 \(\mathrm{V}^{\lambda}\) 与 \(\mathrm{V}^{w\lambda}\) 具有相同的维数。\(\mathrm{V}\) 的权集在 \(\mathrm{W}\) 下稳定。

事实上，\(w\) 具有 \(s^*\) 的形式，其中 \(s \in \mathrm{Aut}_e(\mathfrak{g}, \mathfrak{h})\)（§2，第 2 节，定理 2 的推论）。

备注 2. 根据命题 2 的推论 1 以及 §5，第 3 节的备注 2，谈论 V 关于 g 的典则嘉当子代数的权及其重数是有意义的。

备注 3. 即使不假设 g 可分裂，引理 1 (i) 和命题 2 仍以同样的证明成立。

